\documentclass[11pt]{article}
\usepackage[margin=1in]{geometry}
\usepackage{amsmath,amssymb,amsfonts,mathtools}
\usepackage{array,booktabs,enumitem,graphicx,xcolor,hyperref}
\usepackage[T1]{fontenc}
\usepackage[utf8]{inputenc}
\title{Towards a Classification of Fusion Rule Algebras in Rational Conformal
  Field Theories}
\author{Source-backed arXiv sample}
\date{}
\begin{document}
\maketitle
\section{Extracted Lists}

The list environments below are extracted from arXiv LaTeX source and wrapped for pdfLaTeX validation.

\begin{itemize}
\item In the first one, quantum groups are proposed as the underlying
      algebraic structure of RCFT's~\cite{qgroups}, and this approach is
      particularly successfull in the case of WZW models, although
      its extension to general RCFT's is still unclear. In particular this
      line of thougth, when applied toward any classification program,
      shows severe limitations due to the fact
      that the classification of quantum groups themselves is still an
      open problem.
\item The second approach exploits the relationships between
      RCFT's and three-dimensional Chern-Simons theories~\cite{c-s};
      in particular the
      Hilbert space associated to a constant time slice with charges in the
      three dimensional theories is related to the space of conformal blocks
      of corresponding RCFT's.
\end{itemize}

\begin{itemize}
\item[{\bf P1}:]
     $\cF$ {\em is commutative}: $N_{ij}^k=N_{ji}^k$. This follows from the
     isomorphism (called $\Omega$ in~\cite{ms}) between $V_{ij}^k$ and
     $V_{ji}^k$.
\item[{\bf P2}:]
     $\cF$ {\em is associative}:
     \[
     \sum_m N_{ij}^m N_{km}^l = \sum_m N_{ik}^m N_{jm}^l
     \]
     This follows from existence of the {\em Fusion matrix} isomorphism
     \[
     F\left[\begin{array}{cc}i&k\\j&l\end{array}\right]~:~
     \bigoplus_m V_{im}^j \otimes V_{kl}^m \cong
     \bigoplus_m V_{ml}^j \otimes V_{ik}^m
     \]
     which is related to the assumption of duality.
\item[{\bf P3}:]
     {\em Identity}: there exists in $\Xi$ a vector $\phi_0$ called the
     identity such that for all $i,j\in X$: $N_{0i}^j=\delta_i^j$. The vector
     $\phi_0$ is the one in 1 to 1 correspondence with the HWR $\cH_0$.
\item[{\bf P4}:]
     {\em Charge conjugation}: there exists a one to one map of $X$, say
     $C:i\mapsto \ih$, such that $\hat{\ih}=i$ and $\cF$ is invariant under
     $C$, i.e. $N_{ij}^k=N_{\ih\jh}^{\hat{k}}$ and $N_{ij}^0=\delta_{i\jh}$.
     The matrix $C_{ij}\equiv N_{ij}^0$ uppers and lowers indices in
     $\cF$, and from {\bf P1,P2}, $N_{ijk}$ is totally symmetric in its
     indices.
\item[{\bf P5}:]
     {\em Integrality of structure constants}: $N_{ij}^k\in \bN$. This
     follows from the trivial fact that the dimension of a vector space
     like $V_{ij}^k$ is a non-negative integer.
\end{itemize}
\end{document}
